\begin{figure}[t]
\centering
\pgfplotstableset{col sep=comma}
\begin{tikzpicture}
\begin{axis}[
  font=\small,
  ybar, bar width=5.5pt,
  width=\textwidth, height=0.52\textwidth,
  ymin=0, ymax=10.5, ytick={0,3,6,9},
  ylabel={fix rounds},
  xmin=-0.7, xmax=21.7,
  xtick={0,...,21},
  xticklabels from table={data/review-rounds.csv}{id},
  x tick label style={rotate=60,anchor=east,font=\scriptsize},
  axis y line*=left, axis x line*=bottom,
]
\addplot[fill=okorange, draw=black] table[x expr=\coordindex,y=fix_rounds,col sep=comma] {data/review-rounds.csv};
\draw[okverm, dashed, thick] (axis cs:-0.7,\reviewRoundsMedian) -- (axis cs:21.7,\reviewRoundsMedian);
\node[anchor=west, font=\scriptsize, fill=white, inner sep=1pt] at (axis cs:0.6,3.6) {median \reviewRoundsMedian{} rounds};
\end{axis}
\end{tikzpicture}
\caption{Every package needed at least one fix round in independent review: the median is \reviewRoundsMedian{} and the maximum (\texttt{pelican}) is \reviewRoundsMax{}. One round can fix several defects, so the count is a proxy for review effort, not a defect count. Data: \texttt{paper/data/review-rounds.csv}.}
\label{fig:rounds}
\end{figure}
